% Source: source/普通高考/2014/2014陕西理.pdf, page 1, question 4.
% Flowchart topology: 开始 -> 输入 N -> S=1,i=1 -> a_i=2*S -> S=a_i
% -> i=i+1 -> i>N?; 否 loops to the entry above a_i, 是 -> 输出 -> 结束.
\begin{tikzpicture}[
  x=1cm, y=0.84cm, >=Stealth,
  line width=0.45pt, line cap=round, line join=round,
  every node/.style={anchor=center, align=center,
    execute at begin node={\setlength{\parindent}{0pt}\noindent}},
  flowbox/.style={draw, minimum height=0.70cm, inner xsep=4pt, inner ysep=2pt,
    align=center, execute at begin node={\setlength{\parindent}{0pt}\noindent}},
  terminal/.style={flowbox, rounded corners=0.16cm, minimum width=1.30cm,
    minimum height=0.72cm},
  io/.style={flowbox, trapezium, trapezium left angle=72, trapezium right angle=108,
    minimum width=2.55cm, minimum height=0.78cm},
  process/.style={flowbox, minimum width=2.85cm},
  decision/.style={flowbox, diamond, aspect=3.0, minimum width=3.55cm,
    minimum height=1.00cm, inner sep=1pt},
  branchlabel/.style={anchor=center, align=center, inner sep=0pt,
    execute at begin node={\setlength{\parindent}{0pt}\noindent}}
]
  \node[terminal] (start) at (0,12.00) {开始};
  \node[io] (input) at (0,10.55) {输入 $N$};
  \node[process, minimum width=3.25cm] (init) at (0,9.15) {$S=1,\ i=1$};
  \node[process, minimum width=2.95cm] (double) at (0,7.70) {$a_i=2*S$};
  \node[process, minimum width=2.45cm] (assign) at (0,6.30) {$S=a_i$};
  \node[process, minimum width=2.65cm] (inc) at (0,4.90) {$i=i+1$};
  \node[decision] (test) at (0,3.35) {$i>N$};
  \node[io, minimum width=4.10cm] (output) at (0,1.75) {输出 $a_1,a_2,\cdots,a_N$};
  \node[terminal] (finish) at (0,0.30) {结束};

  \draw[->] (start.south) -- (input.north);
  \draw[->] (input.south) -- (init.north);
  \coordinate (loopJoin) at (0,8.35);
  \draw (init.south) -- (loopJoin);
  \draw[->] (loopJoin) -- (double.north);
  \draw[->] (double.south) -- (assign.north);
  \draw[->] (assign.south) -- (inc.north);
  \draw[->] (inc.south) -- (test.north);
  \draw[->] (test.south) -- node[branchlabel, right=2pt] {是} (output.north);
  \draw[->] (output.south) -- (finish.north);

  \coordinate (loopLeft) at (-3.25,3.35);
  \coordinate (loopTop) at (-3.25,8.35);
  \draw (test.west) -- node[branchlabel, above=2pt] {否} (loopLeft) -- (loopTop);
  \draw[->] (loopTop) -- (loopJoin);
\end{tikzpicture}
